% =====================================================================
%  СЛУЖЕБНОЕ ПРИЛОЖЕНИЕ. Происхождение числовых примеров.
%
%  Назначение (STYLE.md §3): основной текст курса не содержит путей к
%  файлам, имён программ и внутренних обозначений наборов данных —
%  измерение описывается физически. Прослеживаемость нужна для
%  внутренней проверки и живёт здесь.
%
%  При публикации курса вовне: \provenancefalse в preamble.tex —
%  приложение и ключи \srckey исчезают, основной текст не меняется.
% =====================================================================
\section{Происхождение числовых примеров}
\label{sec:provenance}
\markboth{Приложение. Происхождение числовых примеров}{}

\begin{tcolorbox}[enhanced,colback=cPanel,colframe=cGrid,boxrule=0.8pt,arc=2pt]
\small Служебный раздел для внутренней проверки, не входящий в учебное изложение. Каждый ключ
вида \srckey{0.3} в тексте курса соответствует строке таблицы ниже. Пути к файлам даны от корня
репозитория; сокращение \artifact{res/} означает \artifact{research/results/}.
\end{tcolorbox}

\subsection*{Соответствие «описание образца в тексте {$\to$} набор измерений»}

В основном тексте образцы обозначаются физически~--- через параметр плотности. Ниже дано
соответствие внутренним обозначениям наборов данных.

\begin{center}
\footnotesize
\setlength{\tabcolsep}{5pt}
\begin{tabular}{@{}p{0.40\linewidth}p{0.24\linewidth}>{\raggedright\arraybackslash}p{0.30\linewidth}@{}}
\toprule
\textbf{Как назван в курсе} & \textbf{Набор} & \textbf{Примечание} \\
\midrule
плотная решётка, $D/P=0{,}71$ & \artifact{att-11-16-356} & пара поляризаторов; вращается второй \\
повторные установки того же образца, $D/P=0{,}69$ & \artifact{att-11-16-s1,\,-s2,\,-s3} & тот же прибор, другие условия съёмки \\
образец средней плотности, $D/P=0{,}56$ & \artifact{specac} & одиночный поляризатор между эталонными \\
разреженный образец, $D/P=0{,}46$ & \artifact{test\_grid\_40\_20} & одиночный поляризатор \\
наиболее разреженный образец, $D/P=0{,}42$ & \artifact{purewave} & одиночный поляризатор; два дня измерений \\
\bottomrule
\end{tabular}
\end{center}

\begin{tcolorbox}[enhanced,colback=cS2!6,colframe=cS2,boxrule=0pt,leftrule=3pt,arc=1pt]
\small\textbf{Внимание при пересчёте.} Наборы \artifact{att-11-16-*}~--- это ОДИН прибор
(ATT-11-16-CA85, S/N~\#02721) в разных условиях съёмки, а не четыре образца. В курсе они
проходят как повторные установки одного образца; считать их независимыми~--- значит вчетверо
завысить объём выборки.
\end{tcolorbox}

\subsection*{§0. Электродинамика: необходимый минимум}

\begin{center}
\footnotesize
\setlength{\tabcolsep}{5pt}
\begin{tabular}{@{}p{0.07\linewidth}>{\raggedright\arraybackslash}p{0.36\linewidth}>{\raggedright\arraybackslash}p{0.49\linewidth}@{}}
\toprule
\textbf{Ключ} & \textbf{Утверждение в тексте} & \textbf{Источник} \\
\midrule
0.1 & $\tleak\approx0{,}19$~пс из наклона фазы; сдвиг максимума импульса $-0{,}31$~пс &
\artifact{research/SYNTHESIS.md}, линия HN8 ($\dAIC=-579$ поверх HN2);
\artifact{BRIEF\_FOR\_D\_teaching\_ladder.md} §7.1 \\
\addlinespace[2pt]
0.2 & Невязка пятимерного разложения $1{,}4\times10^{-16}$; $A_4/A_2$ $=$ $0{,}2498$ / $0{,}2458$ / $0{,}2549$ / $0{,}2549$ &
\artifact{research/two\_wgp/NOTE\_malus\_linearization.md};
\artifact{res/two\_wgp/a0\_malus\_linearization.json},
\artifact{res/two\_wgp/a0\_malus\_harmonics\_data.json};
\artifact{BRIEF\_FOR\_D\_teaching\_ladder.md} §5.2 \\
\addlinespace[2pt]
0.3 & Период по дифракции $23{,}4\pm0{,}8$~мкм против микрофото $25{,}5\pm0{,}9$~мкм ($1{,}8\sigma$) &
Задача A8 (дифракция лазерной указки), \artifact{res/two\_wgp/a8\_diffraction.json};
сводка~--- \artifact{coordination/BOARD.md}, запись от 2026-08-04 \\
\addlinespace[2pt]
0.4 & Просачивание $\eta\approx0{,}012$ при $D/P=0{,}71$ против $\approx0{,}033$ при $D/P=0{,}46$ &
\artifact{research/SYNTHESIS.md} §2 (модельно-независимая оценка из гармоник) \\
\addlinespace[2pt]
0.5 & $\Deff/\Dphys$ от $0{,}659$ ($D/P=0{,}416$) до $0{,}373$ ($D/P=0{,}710$); точный расчёт $\approx1$ &
\artifact{BRIEF\_FOR\_D\_teaching\_ladder.md} §11.3, колонка ПОЛНОГО углового диапазона.
\textbf{Внимание:} значения с обрезки $(0^\circ,90^\circ)$ смещены и в курсе не используются \\
\addlinespace[2pt]
0.6 & Диаметр перетяжки $1{,}5\ldots2$~мм & \artifact{data\_pool/purewave/README.md}, раздел об
оптической схеме \\
\addlinespace[2pt]
0.7 & Отзыв вывода о связи дефицита $\Deff$ с качеством решётки; угол и участок неразделимы &
Уточнение схемы владельцем 2026-08-04; гипотеза HN15 переведена в статус НЕ ПРОВЕРЕНА,
\artifact{research/hypotheses/HYPOTHESES.md};
корректная постановка~--- \artifact{research/two\_wgp/PROTOCOL\_aperture\_scan.md} \\
\bottomrule
\end{tabular}
\end{center}

\subsubsection*{Рисунки §0}

\begin{center}
\footnotesize
\setlength{\tabcolsep}{5pt}
\begin{tabular}{@{}p{0.20\linewidth}p{0.33\linewidth}>{\raggedright\arraybackslash}p{0.41\linewidth}@{}}
\toprule
\textbf{Рисунок} & \textbf{Что построено} & \textbf{Источник чисел} \\
\midrule
рис.~\ref{fig:harmonics} & тождество для $\cos^4\theta$; измеренные $A_4/A_2$ & тождество~---
чистая тригонометрия; точки~--- ключ~0.2 \\
рис.~\ref{fig:geometry} & два канала решётки, конвенция угла & схема, данные не используются;
конвенция~--- \artifact{research/two\_wgp/model\_2wgp.py} \\
рис.~\ref{fig:scales} & $\lambda=c/\nu$ и полоса периодов & чистая кинематика, данные не
используются \\
рис.~\ref{fig:dp} & шкала $D/P$; $\Deff/\Dphys$ против плотности & ключ~0.5; граница применимости
$D/P<0{,}5$~--- Blanco~A., Fonti~S., Piacente~A. // Infrared Phys. 1986. Vol.~26, №~6.
P.~357--363 \\
\bottomrule
\end{tabular}
\end{center}

\subsection*{Числа, сознательно не использованные}

Эти результаты лежат в материалах проекта, выглядят пригодными и при этом непригодны для учебных
примеров:

\begin{enumerate}[leftmargin=1.6em,itemsep=0.3em]
  \item $\Deff/\Dphys$, полученные на усечённом угловом диапазоне $(0^\circ,90^\circ)$: они
        смещены (для плотного образца $0{,}678$ против $0{,}373$). Берётся только полный
        диапазон.
  \item Значение перекрёстного члена $|\varepsilon_{\rm cross}|=0{,}0345$~--- отозвано:
        соответствует локальному минимуму, хуже глобального на $192{,}5$ по информационному
        критерию. Допустимо только как иллюстрация ошибки многостартовой оптимизации, с явным
        указанием на это (§7).
  \item Утверждение о связи дефицита эффективного диаметра с нерегулярностью решётки~--- не
        проверено (ключ~0.7).
  \item Четыре набора \artifact{att-11-16-*} как «четыре образца»~--- см. врезку выше.
\end{enumerate}
